\subsection{Contesto}
Il contesto in cui si inserisce il progetto \textit{Milan RailSim} è quello della simulazione della rete ferroviaria lombarda, in particolare la rete suburbana e il passante milanese. \\
L'obiettivo del progetto è la simulazione del traffico ferroviario, del comportamento dei treni e della gestione delle stazioni, al fine di analizzare l'efficienza del sistema e identificare eventuali criticità.\\
L'obiettivo principale è quello di analizare le reti suburbane e cercare di analizzare se sia possibile avere un comportamente più simile ad una metropolitana piuttosto che ad un sistema ferroviario, dunque simulare se fosse possibile aumentare il numero di corse e ridurre l'attesa, senza modificare l'infrastruttura già esistente.\\
Dato che sulla stessa infrastruttura circolano, insieme ai suburbani, anche i treni regionali e treni di altre società, non è possibile simulare solamente il comportamento dei suburbani, in quanto non sarebbe corretto: è stata quindi simulata l'intera rete regionale di Trenord. I treni delle altre società, come quelli a lunga percorrenza e ad alta velocità, usano in parte gli stessi binari ma non sono rappresentati nel modello.\\
Il punto più delicato della rete è il Passante ferroviario, una galleria a due binari che attraversa la città e nella quale confluiscono sei linee suburbane. Nell'orario in vigore vi transitano 12 treni all'ora per verso, con intervalli irregolari che vanno da 2 a oltre 9 minuti.\\
Un altro problema che è stato necessario modellare è la presenza di binari unici, dunque tratti di binario utilizzate in entrambi i versi, dunque per mantenere realistico il modello è stato necessario modellare anche questo comportamento.\\

\subsection{Motivazione}
\'E stato scelto di simulare questo caso d'uso, sia per un interesse personale, in quanto da pendolare vivo tutti i giorni le problematiche dei treni in Lombardia, sia per un interesse accademico, in quanto la rete ferroviaria lombarda è un esempio di sistema complesso, in cui il comportamento dell'intero sistema non si ricava dalle regole dei singoli treni, ma dal comportamento collettivo dei treni e dalle interazioni tra di essi.\\

\subsection{Domanda di ricerca e obiettivi}
La domanda da cui parte il lavoro è la seguente: di quanto si può aumentare la frequenza dei treni prima che i ritardi si propaghino a tutta la rete, e quali punti dell'infrastruttura fissano questo limite?

L'obiettivo non è trovare l'orario migliore.
L'obiettivo è comprendere come il sistema reale si comporta e come potrebbe comportarsi in scenari alternativi.
Gli obiettivi da raggiungere sono quattro:
\begin{itemize}
	\item costruire un modello di base utilizzando la rete reale e gli orari ufficiali del servizio, cercando di modellare un sistema realistico
	\item implementazione di uno scenario in cui la rete suburbana ha un aumento delle corse mentre il resto della rete rimane invariato;
	\item implementazione di uno scenario più irrealistico dei primi due, in cui tutta la rete subisce un aumento delle corse e di riduzione dei tempi di attesa, con l'obiettivo di ricercare i colli di bottiglia e dove la rete si satura più velocemente;
	\item per ogni scenario, stimare i costi i costi di servizio, come il costo per la trazione, del personale, ecc...
\end{itemize}

\subsection{Che cosa è stato realizzato}

La simulazione si basa su \textit{railsim}~\cite{kaddoura2024railsim}, un'estensione ferroviaria della piattaforma di simulazione ad agenti MATSim. I dettagli di ogni parte verranno specificati nelle sezioni dedicate.\\

Da railsim sono stati usati:
\begin{itemize}
	\item il motore di simulazione, che fa avanzare i treni sulla rete secondo per secondo;
	\item il moto dei treni, con accelerazione, frenata e velocità massima proprie di ogni tipo di treno;
	\item l'occupazione dei binari: un treno entra in una tratta solo se è libera, altrimenti attende;
	\item la rappresentazione di una tratta a binario unico come risorsa condivisa dai due versi di marcia;
	\item la possibilità di deviare un treno su un altro binario della stessa stazione;
	\item i punti di estensione, attraverso i quali sono state inserite le regole aggiuntive.
\end{itemize}

Per il progetto sono stati sviluppati:
\begin{itemize}
	\item la costruzione della rete a partire dall'orario pubblico di Trenord e dai dati di OpenStreetMap;
	\item le stazioni principali, con specifica di dettaglio a livello di binario, con gli ingressi, le uscite e i binari di ricovero;
	\item il caricamento delle corse dall'orario pubblico e la loro assegnazione ai treni e ai binari, informazioni che l'orario non contiene;
	\item le regole per gli incroci sul binario unico, che stabiliscono quando un treno può entrare in una tratta e dove attende il treno opposto;
	\item la misura della puntualità, il calcolo del consumo di energia e il calcolo dei costi, che railsim non prevede;
	\item i generatori degli scenari, che aggiungono corse all'orario reale;
	\item un'applicazione grafica che mostra i treni sulla mappa durante la simulazione, archivia ogni simulazione e ne presenta i risultati in tabelle e grafici.
\end{itemize}

\subsection{Metodo di lavoro}
La realizzazione del progetto è stata guidata dai dati ottenuti dalla simulazioni e dal confronto con i dati reali.\\
Lo sviluppo è stato TDD (Test Driven Development), ovvero prima della scrittura del codice sono state scritte le suite di test e i test unitari che doveno passare una volta implementata la funzionalità.\\
Al termine dello sviluppo veniva eseguito lo scenario di simulazione su una giornata completa, per verificare se il comportamento era corretto o si verificano situazioni anomale.\\
Per lo sviluppo del codice e l'esecuzione dei test è stato è stato d'aiuto \textit{Claude Code}, l'assistente di programmazione di Anthropic, con i modelli Opus 5.5 e Fable 5.1.